\chapter{A estrutura da automação}
\label{cap:estrutura}

\begin{objetivos}
Ao final deste capítulo você deve ser capaz de:
\begin{itemize}
  \item descrever os níveis hierárquicos de um sistema de automação industrial;
  \item dizer o que trafega em cada nível e em que direção;
  \item relacionar cada nível às tecnologias e protocolos usualmente empregados;
  \item localizar o CLP nessa hierarquia e explicar por que ele não está sozinho.
\end{itemize}
\end{objetivos}

\section{Do chão de fábrica à gestão}

Um sistema de automação não é apenas um controlador ligado a sensores e atuadores/elementos finais
de controle. É uma hierarquia em que cada nível consome informação do nível inferior e devolve
decisão ao nível de baixo. Essa organização, chamada de \emph{pirâmide da automação}, existe porque
os requisitos mudam a cada degrau: no campo importa o tempo de resposta; na gestão, importa o
histórico e o indicador do mês.

\begin{figure}[H]
  \centering
  \includegraphics[width=\linewidth]{03-piramide-niveis.jpeg}
  \caption{Níveis da automação industrial e fluxo de informação entre eles.}
  \label{fig:piramide}
\end{figure}

O fluxo tem duas direções (Figura~\ref{fig:piramide}): os dados sobem --- medida, estado, evento ---
e os comandos descem --- meta, \emph{setpoint}, ordem de produção.

\index{pirâmide da automação}\index{automação industrial!níveis}

\section{Os cinco níveis}

\begin{table}[H]
  \centering
  \caption{Níveis da automação: função, equipamentos e tecnologias típicas.}
  \label{tab:niveis}
  \begin{tabular}{p{2.6cm} p{4.3cm} p{3.4cm} p{4.3cm}}
    \toprule
    Nível & Função & Equipamentos & Tecnologias usuais \\
    \midrule
    5. Corporativo
      & Gestão financeira, administrativa e de recursos; indicadores de resultado
      & ERP, sistemas de \emph{business intelligence}
      & Redes corporativas, bancos de dados; acesso raramente direto aos níveis fabris \\
    4. Gerencial
      & Planejamento detalhado da produção, logística, indicadores e apoio à manutenção preditiva
      & MES, PIMS
      & Ethernet de alto desempenho, banco de dados histórico \\
    3. Supervisão
      & Operação assistida, alarmes, telas, histórico e relatórios do processo
      & SCADA, IHM, \emph{historian}, servidores de dados
      & Ethernet TCP/IP; OPC como camada de interoperabilidade entre sistemas \\
    2. Controle
      & Executar o controle automático: lógica, intertravamento, malhas
      & CLP, SDCD, remotas de I/O
      & Profibus DP, DeviceNet, ControlNet, Modbus TCP, EtherNet/IP \\
    1. Campo
      & Transformar grandeza física em sinal e executar a ação sobre o processo
      & Sensores, transmissores, atuadores, posicionadores
      & Redes de campo com baixa largura de banda: CAN, AS-i, Modbus, Profibus PA, Fieldbus Foundation H1 \\
    \bottomrule
  \end{tabular}
\end{table}

\paragraph{Nível 1 --- Campo.} É onde o processo encontra a instrumentação. Os dispositivos podem ser
ponto a ponto, ligados diretamente ao cartão de entrada e saída, ou microprocessados e inteligentes,
gerando volume grande de dados. As redes aqui são heterogêneas e, em geral, de baixa largura de
banda.

\paragraph{Nível 2 --- Controle.} É onde vive o CLP. Exige inteligência local suficiente para
executar lógica, intertravamento e malhas de controle, e comunicação mais sofisticada que a do
campo: protocolos com velocidade maior e capacidade de resposta em tempo real.

\paragraph{Nível 3 --- Supervisão.} Concentra a operação. Aqui ficam as telas, o banco de dados do
processo, o registro de eventos e alarmes. Neste nível é comum encontrar conversores de protocolo
usando OPC, exatamente para garantir interoperabilidade entre sistemas de fabricantes diferentes
(Capítulo~\ref{cap:interoperabilidade}).

\paragraph{Nível 4 --- Gerencial.} Transforma dados em informação de planejamento. Ferramentas como
MES e PIMS produzem relatórios mais elaborados que os da supervisão, porque trabalham com massa de
dados maior e com histórico.

\paragraph{Nível 5 --- Corporativo.} Trata da gestão de recursos da empresa. Trabalha com resultado
financeiro e, tradicionalmente, tem pouco acesso direto aos dados técnicos dos níveis inferiores ---
é justamente essa fronteira que as arquiteturas de IIoT vieram discutir
(Capítulo~\ref{cap:iiot}).

\section{O que mudou na pirâmide}

A pirâmide clássica pressupõe um fluxo lento de baixo para cima e uma separação clara entre rede
fabril e rede administrativa. Dois movimentos alteraram esse desenho:

\begin{enumerate}
  \item \textbf{Redes de campo substituíram a fiação ponto a ponto}, reduzindo drasticamente o
        cabeamento entre o processo e o sistema de controle. Em vez de levar todo sinal até o rack
        do CLP, os sinais vão para remotas de I/O instaladas no campo
        (Capítulo~\ref{cap:redes-scada}).
  \item \textbf{A camada de informação se sobrepõe aos níveis}, e não os substitui: dados de um
        sensor de campo passam a interessar diretamente à manutenção e à gestão, sem atravessar
        necessariamente todos os degraus. A governança desse acesso é um problema de arquitetura e
        de segurança, não apenas de rede.
\end{enumerate}

\begin{atencao}
Nível hierárquico não é o mesmo que prioridade de tempo. O nível de campo e o de controle respondem
em milissegundos; o gerencial trabalha em horas ou dias. Confundir essas escalas é a origem de
vários projetos mal dimensionados --- e de acidentes, quando se tenta fazer controle rápido por
caminho lento.
\end{atencao}

\begin{pratica}
Visite uma planta ou examine a arquitetura de um sistema que você conheça e preencha a
Tabela~\ref{tab:niveis} com os equipamentos \emph{reais} de cada nível. Depois responda:
(a) algum nível está ausente ou foi fundido com outro?
(b) qual é o caminho que um dado de sensor percorre até virar indicador de gestão?
(c) existe alguma ligação direta entre os níveis 1--2 e o nível 4 ou 5? Quem autorizou?
\end{pratica}

\begin{exercicios}
\begin{enumerate}[label=\alph*)]
  \item Explique por que a largura de banda exigida no nível de campo é menor que a exigida no
        nível de supervisão.
  \item Um instrumento inteligente gera 40 parâmetros de diagnóstico, mas o projeto só usa a variável
        de processo. Cite duas consequências práticas e diga para qual nível esses parâmetros
        interessam.
  \item Descreva o caminho de um alarme desde a sua detecção no CLP até o registro no nível
        gerencial, nomeando cada equipamento envolvido.
  \item Por que o uso de OPC aparece tipicamente no nível de supervisão e não no nível de campo?
  \item A direção ``gerência $\rightarrow$ chão de fábrica'' leva que tipo de informação? Dê dois
        exemplos.
\end{enumerate}
\end{exercicios}
